\chapter{Questões de Pesquisa}
\label{chap:questoes}

Estabelecido no capítulo anterior que a eficácia de uma estratégia de
retreinamento depende conjuntamente do tipo e da severidade do drift
enfrentado e do classificador monitorado, e que essa dependência raramente é
investigada de forma sistemática, este capítulo formaliza as questões que
orientam a investigação empírica deste trabalho.

\section{Questão Principal}

\begin{quote}
\textit{Como a eficácia de diferentes estratégias de retreinamento varia em
função do tipo e da severidade do drift enfrentado e do classificador
monitorado, em ambiente de aprendizado em stream?}
\end{quote}

Essa questão se desdobra em três questões específicas, cada uma associada a
uma das variáveis manipuladas no experimento fatorial.

\textbf{QE1 — Tipo de drift.} A eficácia relativa das estratégias de
retreinamento varia entre os padrões temporais de drift abrupto, gradual,
incremental e recorrente \cite{gama2014,lu2018}? Espera-se, em particular,
que a reativação de modelos históricos apresente vantagem específica sob
drift recorrente, na medida em que descartar um modelo que voltará a ser
válido representa custo evitável \cite{sun2018,li2021}.

\textbf{QE2 — Severidade do drift.} A eficácia relativa das estratégias de
retreinamento varia conforme a magnitude da mudança, independentemente do seu
padrão temporal \cite{wang2022elastic,wang2024quadcdd}? Espera-se que
estratégias mais radicais, como o reset completo, apresentem vantagem sob
alta severidade, e estratégias mais conservadoras, como a janela deslizante
ou o retreinamento incremental, sob baixa severidade.

\textbf{QE3 — Classificador monitorado.} A ordenação de eficácia entre as
estratégias de retreinamento se mantém estável entre classificadores com
mecânicas distintas, ou a estratégia mais eficaz depende do classificador
monitorado \cite{webb2016,wang2022dynamic}?

\section{Questão Derivada}

Como desdobramento exploratório, e de caráter descritivo antes que causal,
este trabalho investiga ainda:

\begin{quote}
\textit{É possível relacionar sinais observáveis do detector adotado — como
a magnitude da divergência calculada, utilizada como proxy de severidade, a
velocidade com que essa divergência varia, indicativa do padrão temporal, e
a semelhança entre a distribuição atual e distribuições históricas, indício
de recorrência — aos padrões de eficácia identificados no experimento
fatorial, de modo a esboçar uma orientação preliminar sobre qual estratégia
de retreinamento adotar diante de cada condição observável?}
\end{quote}

Essa questão não pressupõe uma resposta definitiva, mas busca aproximar os
resultados obtidos de uma forma de orientação aplicável a pipelines de
MLOps, nos quais o monitoramento contínuo e os mecanismos de realimentação
já são reconhecidos como componentes de primeira ordem, ainda carentes de
critérios explícitos para a escolha da estratégia de resposta acionada após
a detecção \cite{kreuzberger2023}.